\documentclass[11pt,a4paper]{article}
\usepackage[margin=2.5cm]{geometry}
\usepackage{amsmath,amssymb}

\newcommand{\R}{\mathbb{R}}
\newcommand{\norm}[1]{\left\| #1 \right\|}
\newcommand{\xs}{x^\star}

\title{Homework 12.10.26 -- Exercises 9 and 10}
\author{Gor Melkumyan}
\date{}

\begin{document}
\maketitle

\section*{Exercise 9}

Given $\mu I \preceq \nabla^2 f \preceq L I$, $0 < \mu \le L$, and
$T_\gamma(x) := x - \gamma \nabla f(x)$.

\subsection*{1)}
\[
  T_\gamma(x) - T_\gamma(y) = \int_0^1 \nabla T_\gamma\bigl(y + t(x-y)\bigr)(x-y)\,dt
  \;\Rightarrow
\]
\begin{align*}
  \norm{T_\gamma(x) - T_\gamma(y)}
  &= \norm{\int_0^1 \nabla T_\gamma\bigl(y + t(x-y)\bigr)(x-y)\,dt}
  \le \int_0^1 \norm{\nabla T_\gamma\bigl(y + t(x-y)\bigr)}\norm{x-y}\,dt \\
  &\le \int_0^1 \sup_{z \in \R^d}\norm{\nabla T_\gamma(z)}\,dt \cdot \norm{x-y}
  = \sup_{z \in \R^d}\norm{I - \gamma \nabla^2 f(z)}\,\norm{x-y}.
\end{align*}
Since $\mu I \preceq \nabla^2 f \preceq L I$, the eigenvalues $\lambda_i$ of the
symmetric matrix $\nabla^2 f(z)$ lie in $[\mu, L]$, so
\[
  \norm{I - \gamma \nabla^2 f(z)} = \max_i |1 - \gamma\lambda_i|
  \le \max\{|1-\gamma L|, |1-\mu\gamma|\}
\]
($|1-\gamma\lambda|$ is convex in $\lambda$, so its maximum on $[\mu, L]$ is at an endpoint)
$\Rightarrow$
\[
  \norm{T_\gamma(x) - T_\gamma(y)} \le \norm{x-y}\cdot\max\{|1-\gamma L|, |1-\mu\gamma|\}.
\]

\subsection*{2)}
If $\gamma = \frac{2}{\mu+L}$ $\Rightarrow$
$|1-\gamma L| = |1-\mu\gamma| = \frac{L-\mu}{L+\mu}$.

If $\gamma > \frac{2}{\mu+L}$ $\Rightarrow$ $L\gamma > \frac{2L}{\mu+L}$ $\Rightarrow$
$1 - L\gamma < 1 - \frac{2L}{\mu+L} = \frac{\mu-L}{\mu+L} \le 0$ $\Rightarrow$
$|1-L\gamma| > \frac{L-\mu}{L+\mu}$ $\Rightarrow$
$\max\{|1-\gamma L|, |1-\mu\gamma|\} > \frac{L-\mu}{L+\mu}$.

If $\gamma < \frac{2}{\mu+L}$ $\Rightarrow$ $\mu\gamma < \frac{2\mu}{\mu+L}$ $\Rightarrow$
$1 - \mu\gamma > 1 - \frac{2\mu}{\mu+L} = \frac{L-\mu}{L+\mu} \ge 0$ $\Rightarrow$
$|1-\mu\gamma| > \frac{L-\mu}{L+\mu}$ $\Rightarrow$
$\max\{|1-\gamma L|, |1-\mu\gamma|\} > \frac{L-\mu}{L+\mu}$.

$\Rightarrow$ $\gamma = \frac{2}{\mu+L}$ is optimal.

\subsection*{3)}
Since $\nabla f(\xs) = 0$, we have $T_\gamma(\xs) = \xs$, so with $\gamma = \frac{2}{\mu+L}$
\[
  \norm{x^{k+1} - \xs} = \norm{T_\gamma(x^k) - T_\gamma(\xs)}
  \le \frac{L-\mu}{L+\mu}\norm{x^k - \xs}
  \;\Rightarrow\;
  \norm{x^k - \xs} \le \norm{x^0 - \xs}\cdot\left(\frac{L-\mu}{L+\mu}\right)^k .
\]
In Theorem 15, $\norm{x^k - \xs}^2 \le (1-\gamma\mu)^k\norm{x^0 - \xs}^2$, where
$\gamma \le \frac{1}{L}$; the best choice is $\gamma = \frac{1}{L}$
\[
  \Rightarrow\; \norm{x^k - \xs}^2 \le \left(1 - \frac{\mu}{L}\right)^k\norm{x^0 - \xs}^2
  = \left(\frac{L-\mu}{L}\right)^k\norm{x^0 - \xs}^2 .
\]
Squaring our bound, $\norm{x^k - \xs}^2 \le \left(\frac{L-\mu}{L+\mu}\right)^{2k}\norm{x^0 - \xs}^2$, and
\[
  \frac{L-\mu}{L} : \left(\frac{L-\mu}{L+\mu}\right)^2
  = \frac{(L+\mu)^2}{L(L-\mu)} > 1
  \quad (\mu < L),
\]
so the rate with $\gamma = \frac{2}{\mu+L}$ is better, by the factor
$\frac{(L+\mu)^2}{L(L-\mu)}$ per iteration.

\section*{Exercise 10}

If $f$ is $L$-smooth and $\mu$-strongly convex, then by (18) and Exercise 8
\[
  \frac{1}{L}\norm{\nabla f(x^k) - \nabla f(\xs)}^2
  \le 2D_f(x^k, \xs)
  \le \frac{1}{\mu}\norm{\nabla f(x^k) - \nabla f(\xs)}^2 .
\]
Dividing by $\norm{\nabla f(x^k) - \nabla f(\xs)}^2$, (i) and (ii) are true:
$\frac{1}{L} \le \gamma_k \le \frac{1}{\mu}$.

For (iii): by (i), $1 - \gamma_k L \le 0$, so the last term of Lemma 14 cannot be
dropped. Instead we stop the proof of Lemma 14 before using (25), with
$r^k = x^k - \xs$:
\[
  \norm{r^{k+1}}^2 \le (1 - \gamma_k\mu)\norm{r^k}^2 - 2\gamma_k D_f(x^k, \xs)
  + \gamma_k^2\norm{\nabla f(x^k) - \nabla f(\xs)}^2 .
\]
By the definition of $\gamma_k$,
$\gamma_k^2\norm{\nabla f(x^k) - \nabla f(\xs)}^2 = 2\gamma_k D_f(x^k, \xs)$
$\Rightarrow$
\[
  \norm{x^{k+1} - \xs}^2 \le (1 - \gamma_k\mu)\norm{x^k - \xs}^2 .
\]

\end{document}
